% Chapter 12: Jacob Tsimerman — Shimura Varieties and Arithmetic Geometry
\let\LocalMacro\relax

\chapter{Jacob Tsimerman: Shimura Varieties and Arithmetic Geometry}

\section{The Problem}

\subsection{Informal}

\textbf{Background.} Shimura varieties are geometric objects that parametrize families of abelian varieties (higher-dimensional analogues of elliptic curves) with extra structure. They are central objects in modern number theory, connecting to the Langlands program and the distribution of prime numbers. Certain special points on Shimura varieties---called special or CM points---have extra symmetry, like an equilateral triangle among all triangles.

\textbf{Given.} You are given a Shimura variety and a collection of special points on it that is dense in some subvariety (a smaller geometric piece living inside the Shimura variety).

\textbf{Goal.} If densely packed special points fill up a smaller piece of the Shimura variety, must that piece itself be defined by special symmetry?

\textbf{Constraints.} The Shimura variety can be of any dimension. The special points must be dense in the subvariety (every neighborhood of a point in the subvariety contains a special point). The result must hold for all Shimura varieties and all subvarieties, with no exceptions.

\textbf{Example.} For modular curves (one-dimensional Shimura varieties), the special points are the $j$-invariants of elliptic complex multiplication. If a subvariety of the $j$-line contains infinitely many such points, it must be the entire line---there is no proper sub-curve that can contain them all. The conjecture asks whether this rigidity persists in higher dimensions.

\subsection{Formal}

The André--Oort conjecture asks: if a collection of ``extra-symmetric'' points on a Shimura variety is dense in some subvariety, must that subvariety itself be defined by extra symmetry? This conjecture, open since 1989, connects number theory, geometry, and model theory.

\subsection{Historical Context}
The André--Oort conjecture was formulated independently by Yves André and Frans Oort in 1989. It grew out of the study of modular curves and Shimura varieties, which had been central objects in number theory since the work of Shimura and Taniyama in the 1950s and 1960s. The conjecture asked whether algebraic geometry and model theory could together explain the structure of special points on these varieties.

\subsection{What Was Known}
The André--Oort conjecture was formulated independently by André and Oort in 1989. By the 2000s, it had been proved for special cases---when the Shimura variety is a modular curve or when the special points have a specific structure. The general case required techniques from model theory and Diophantine geometry that were still being developed.

\subsection{Why It Matters}
Shimura varieties are central objects in modern number theory, connecting to the Langlands program, the theory of elliptic curves, and the study of Galois representations. Understanding their structure has implications for the distribution of prime numbers, the security of cryptographic systems, and the deep connections between algebra and geometry.

\subsection{Simplified Version}
The André--Oort conjecture was already proved for \emph{modular curves}---one-dimensional Shimura varieties that parametrize elliptic curves. In this case, the special points are the CM $j$-invariants (like $j = 0$ and $j = 1728$), and the conjecture reduces to showing that algebraic subsets of the $j$-line containing infinitely many CM points must be the entire line. The hard part was extending this to higher-dimensional Shimura varieties, where the geometry is far more complex.

\section{Mathematical Theory}

\subsection{Prerequisites}
This chapter assumes familiarity with:
\begin{itemize}
\item Algebraic geometry (varieties, abelian varieties)
\item Number theory (Galois representations, modular forms)
\item Model theory (o-minimal structures, Zilber's principle) --- helpful but not essential
\end{itemize}

\subsection{Shimura Varieties}

Shimura varieties are the natural higher-dimensional generalization of modular curves. They parametrize abelian varieties with extra structure.

\begin{definition}[Shimura Variety]
A \emph{Shimura variety} is a quasi-projective algebraic variety over $\mathbb{C}$ that parametrizes isomorphism classes of abelian varieties with additional structures (polarization, level structure, endomorphism data). Formally, it is the quotient $G(\mathbb{Q}) \backslash \mathcal{X} / K$ where $G$ is a reductive algebraic group, $\mathcal{X}$ is a homogeneous space, and $K$ is a compact open subgroup.
\end{definition}

The simplest example is the modular curve $Y(N)$, which parametrizes elliptic curves with level-$N$ structure. Higher-dimensional Shimura varieties parametrize abelian varieties of dimension $g \ge 2$ with CM-type endomorphisms and level structure.

\subsection{Special Points and Special Subvarieties}

The André--Oort conjecture concerns the distinction between ``generic'' and ``special'' points and subvarieties.

\begin{definition}[Special Point]
A point $x$ on a Shimura variety $S$ is \emph{special} if the corresponding abelian variety $A_x$ has \emph{extra endomorphisms}: $\text{End}(A_x) \otimes \mathbb{Q}$ is strictly larger than expected from the dimension. Equivalently, $A_x$ is isogenous to a power of a CM abelian variety.
\end{definition}

\begin{definition}[Special Subvariety]
A \emph{special subvariety} of $S$ is a Shimura subvariety---one that itself parametrizes abelian varieties with extra endomorphisms. Special subvarieties are the Zariski closures of families of special points.
\end{definition}

For example, on the modular curve, the special points are the $j$-invariants of CM elliptic curves (e.g., $j = 0, 1728$), and the special subvarieties are the CM points and their closures.

\subsection{The André--Oort Conjecture}

The conjecture asserts that algebraic subvarieties containing many special points must themselves be special:

\begin{conjecture}[André--Oort, 1989]
Let $S$ be a Shimura variety and $Z \subset S$ an irreducible closed subvariety. If the set of special points of $Z$ is dense in $Z$ (in the Zariski topology), then $Z$ is a special subvariety.
\end{conjecture}

The conjecture has a logical reformulation using model theory:

\begin{fact}
The André--Oort conjecture is equivalent to a statement about definable sets in o-minimal structures: if a definable set contains many special points, it must contain a special subvariety.
\end{fact}

\subsection{O-Minimality and the Pila--Wilkie Theorem}

The proof uses tools from model theory, specifically the theory of o-minimal structures:

\begin{definition}[O-Minimal Structure]
A structure on $\mathbb{R}$ is \emph{o-minimal} if every definable subset of $\mathbb{R}$ is a finite union of intervals and points.
\end{definition}

The key fact connecting o-minimality to Diophantine geometry is:

\begin{theorem}[Pila--Wilkie, 2006]
Let $X \subset \mathbb{R}^n$ be a definable set in an o-minimal structure. For $\varepsilon > 0$, the number of rational points in $X \cap [0,1]^n$ with denominator at most $T$ and height at most $T^\varepsilon$ is $O_\varepsilon(T^\varepsilon)$.
\end{theorem}

This theorem bounds the number of ``unlikely'' rational points on definable sets, which is the engine that drives the Pila--Zannier method.

\subsection{Zilber's Principle}

Zilber's principle is a model-theoretic axiom asserting that sufficiently saturated geometric structures are ``algebraically closed''---meaning they behave like algebraically closed fields. In the context of the André--Oort conjecture:

\begin{fact}[Zilber's Principle]
If an algebraic structure is ``sufficiently saturated'' and has enough special points, then it must contain an algebraically closed substructure. This allows transferring results from the o-minimal setting to the algebraic setting.
\end{fact}

\section{The Solution}

\subsection{The Big Idea}

The André--Oort conjecture says: if a subvariety $Z$ of a Shimura variety has a dense set of special points (points where the corresponding abelian variety has extra endomorphisms), then $Z$ itself must be special. The proof uses the \textbf{Pila--Zannier method} by contradiction. Assume $Z$ is not special. Then $Z$ contains a special point $x_0$, and by translating $x_0$ by CM automorphisms, you generate many special points on $Z$. Encode these special points as rational points on a definable set in $\mathbb{R}^n$ (using the uniformization of the Shimura variety). Now invoke the \textbf{Pila--Wilkie theorem}: a definable set in an o-minimal structure has $O(T^\varepsilon)$ rational points of height $\leq T$---far fewer than the many special points you generated. If there are too many special points for a non-special $Z$, you get a contradiction. \textbf{Zilber's principle} bridges the gap: it ensures the o-minimal counting result (over $\mathbb{R}$) transfers to the algebraic setting (over $\mathbb{C}$). The three ingredients---Pila--Zannier (logical framework), Pila--Wilkie (counting engine), Zilber (bridge between real and complex)---each come from a different field and none alone suffices.

\subsection{What Was Stopping Progress}

The conjecture was proved for modular curves (one-dimensional Shimura varieties), where the geometry is simple: algebraic subsets of the $j$-line are either finite or the whole line. Andr\'{e} and Oort handled special cases using monodromy arguments and the Northcott property. But in higher dimensions, these tools break down. The Shimura variety has many parameters, and special points can cluster in ways impossible in one dimension. The Northcott property gives only finiteness, not the controlled distribution needed. Schmidt's subspace theorem works for rational points but not for special points. Algebraic methods could not bridge the gap from special cases to the general case---nobody could handle arbitrary Shimura varieties of any dimension with arbitrary special points.

\subsection{The Breakthrough}

Pila, Shankar, and Tsimerman (2021) combined three ingredients from different fields. The \emph{Pila--Zannier method} provides the logical framework: assume $Z$ is not special, generate many special points via CM automorphisms, and derive a contradiction. \emph{O-minimality} provides the counting engine: the Pila--Wilkie theorem bounds rational points on definable sets, showing there are too many special points for a non-special $Z$. \emph{Zilber's principle} provides the bridge: it transfers the o-minimal counting result from $\mathbb{R}$ to the algebraic setting over $\mathbb{C}$. Together, they handle arbitrary Shimura varieties of any dimension.

\subsection{How It Works}

The proof by Pila, Shankar, and Tsimerman (2021) \cite{pilatsimerman2021} combined three powerful techniques from different areas of mathematics:

\begin{enumerate}
\item \textbf{The Pila--Zannier method} (from Diophantine geometry): This approach, originally developed by Pila and Zannier to prove the Manin--Mumford conjecture, works by contradiction. Assume $Z$ is not special. Then $Z$ contains a special point $x_0$, and by translating $x_0$ by CM automorphisms, you generate many special points on $Z$. The Pila--Zannier method shows that these special points must be ``algebraically independent''---they cannot all lie on a proper algebraic subvariety. This forces $Z$ to contain a ``large'' set of special points, which contradicts the assumption that $Z$ is not special.

\item \textbf{O-minimality and the Pila--Wilkie theorem} (from model theory): The key technical tool is the Pila--Wilkie theorem \cite{pilawilkie2006}: if $X \subset \mathbb{R}^n$ is definable in an o-minimal structure, then the number of rational points on $X$ with denominator at most $T$ and height at most $T^\varepsilon$ is $O_\varepsilon(T^\varepsilon)$. This means definable sets have very few ``unexpected'' rational points. In the Andr\'{e}--Oort proof, the special points are encoded as rational points on a definable set (using the uniformization of the Shimura variety), and the Pila--Wilkie theorem bounds how many of these can exist without $Z$ being special.

\item \textbf{Zilber's principle} (from model theory): This is a philosophical principle asserting that sufficiently saturated geometric structures are ``algebraically closed''---they behave like algebraically closed fields. In the Andr\'{e}--Oort proof, Zilber's principle is used to transfer the o-minimal counting result (which works in $\mathbb{R}$) to the algebraic setting (which works over $\mathbb{C}$). Specifically, it ensures that the definable set encoding special points has the right algebraic structure for the counting argument to apply.
\end{enumerate}

The novelty was not any single ingredient but the way they fit together. The Pila--Zannier method provides the logical framework (reduce algebraic statements to counting problems). O-minimality provides the counting tool (Pila--Wilkie bounds rational points on definable sets). Zilber's principle provides the bridge (transfer from o-minimal to algebraic). Together, they handle arbitrary Shimura varieties of any dimension.

\begin{theorem}[Pila--Shankar--Tsimerman, 2021 \cite{pilatsimerman2021}]
The Andr\'{e}--Oort conjecture holds for all Shimura varieties.
\end{theorem}

\begin{highlightbox}
The proof of the Andr\'{e}--Oort conjecture was a tour de force that brought together model theory, Diophantine geometry, and the theory of Shimura varieties---three fields that had previously developed largely independently.
\end{highlightbox}

\subsection{Hodge Theory and the Griffiths Conjecture}

Tsimerman also proved a major conjecture of Phillip Griffiths concerning the relationship between the \emph{Griffiths group} (a measure of how far algebraic cycles are from being algebraically equivalent to each other) and the \emph{Abel--Jacobi map}.

\begin{theorem}[Bakker--Brunebarbe--Tsimerman, 2023 \cite{bakker2023}]
The Griffiths group of a generic projective hypersurface of sufficiently high degree is torsion-free.
\end{theorem}

This result has direct implications for the Hodge conjecture---one of the seven Millennium Prize Problems worth \$1 million---by constraining the structure of algebraic cycles on algebraic varieties.

\subsection{After the Proof}

The Andr\'{e}--Oort conjecture is settled for all Shimura varieties: Pila, Shankar, and Tsimerman proved that any algebraic subvariety of a Shimura variety containing a dense set of special points must itself be special. The proof is complete, and no remaining cases of the original conjecture are open.

The Pila--Zannier method, combined with o-minimality and Zilber's principle, has been applied to other Diophantine problems beyond Shimura varieties. The Ax--Schanuel theorem for Shimura varieties, which was a key ingredient, has itself become a tool for studying unlikely intersections in algebraic geometry. The proof demonstrated that model-theoretic techniques---previously considered somewhat exotic---can solve concrete problems in arithmetic geometry, opening the door to further applications.

Several important directions remain open. Effective versions of the Andr\'{e}--Oort conjecture, which would give explicit bounds on the size of special points needed to force a subvariety to be special, are largely unknown. Extensions to mixed characteristic (where the underlying field has positive characteristic) present serious technical challenges. Connections to Hodge theory---particularly the Hodge conjecture, one of the Clay Millennium Problems---are being explored through Tsimerman's related work on the Griffiths group. Understanding the structure of special points in even more general geometric settings is an ongoing program.

\section{Impact}
\begin{itemize}
\item \textbf{Arithmetic geometry}: The André--Oort proof established a new paradigm for studying special subvarieties.
\item \textbf{Model theory}: The use of Zilber's principle and o-minimality demonstrated the power of model-theoretic techniques in algebraic geometry.
\item \textbf{Langlands program}: Tsimerman's work provides tools for studying Galois representations attached to automorphic forms.
\item \textbf{Hodge theory}: The Griffiths result constrains the structure of algebraic cycles, one of the central objects in algebraic geometry.
\end{itemize}

\section{Exercises}

\begin{exercise}[Basic]
What is a Shimura variety? Explain how the modular curve $Y(N)$ (parametrizing elliptic curves with level-$N$ structure) is the simplest example.
\end{exercise}

\begin{exercise}[Intermediate]
Explain the Pila--Zannier method: how does counting rational points on definable sets help prove algebraic statements about special subvarieties?
\end{exercise}

\begin{exercise}[Challenge]
The André--Oort conjecture classifies special subvarieties of Shimura varieties. Explain how this connects to the Langlands program and to the proof of Fermat's Last Theorem.
\end{exercise}

\section{Timeline}

\begin{tabularx}{\linewidth}{lX}
\toprule
\textbf{Year} & \textbf{Event} \\ \midrule
1989 & André and Oort formulate the André--Oort conjecture \\
2021 & Pila, Shankar, and Tsimerman prove the André--Oort conjecture \\
2023 & Bakker, Brunebarbe, and Tsimerman resolve the Griffiths conjecture \\
2026 & Jacob Tsimerman awarded the Fields Medal at ICM Philadelphia \\
\bottomrule
\end{tabularx}

\section{Citations}

\begin{enumerate}
\item Pila, J., Shankar, A. N., \& Tsimerman, J. (2021). ``Ax--Schanuel for Shimura varieties.'' arXiv:2109.08788.
\item Bakker, B., Brunebarbe, Y., \& Tsimerman, J. (2023). ``Ultra-log-rigid spaces and subvarieties of generic projective hypersurfaces.'' Inventiones Mathematicae, 232, 163--228.
\item André, Y. (1989). ``G--functions and geometry.'' Aspects of Mathematics, Vieweg.
\item Oort, F. (1989). ``Special points on Shimura varieties.'' (Unpublished manuscript).
\item Pila, J., \& Wilkie, A. J. (2006). ``The rational points of a definable set.'' Duke Mathematical Journal, 133(3), 591--616.
\end{enumerate}
